% Appendix B: the technical detail behind the plain text -- created by loop0015
% item 08x (2026-11-05), the plain-English rewrite. Every passage here was
% moved out of Sections 2, 4, 6 and 7 unchanged in content: every line number
% is at commit 4eaee875 of the clone (read with cat/sed -n only), every
% manuscript quotation gives the page of paper.pdf and the line of
% reports/scratch/eth_audit/paper.txt, and every number names its script in
% the section that cites it. The plain text cites each subsection by label.
\section{Technical detail behind the plain text}
\label{app:tech}
\begingroup\sloppy

The body of the paper is written so that a reader with no background in
complexity theory or in proof assistants can follow it. The passages that
could not be made plain without losing precision are collected here, each
under the heading of the section that cites it. Nothing in this appendix
is new; it is the dense form of what the body says in words.

\subsection{The definitions of the comparator file, line by line}
\label{app:tech-definitions}

This is the detail behind Section~\ref{sec:statement-definitions}. Line
numbers are those of
\clonefile{lean/ComparatorChallenges/SubexponentialSAT.lean} at commit
\texttt{4eaee875}.

\paragraph{Formulas and the count (lines 13--51).}
A \texttt{Literal V} is a variable with a sign (lines 19--22); a
\texttt{Clause V} is a \texttt{Finset} of literals and a \texttt{CNF V} a
\texttt{Finset} of clauses (lines 24--25). An assignment is any function
\(V \to \mathtt{Bool}\) (line 17). A literal evaluates to the Boolean
``the assignment agrees with its sign'' (lines 29--30), a clause to ``some
literal is true'' (lines 32--33), a formula to ``every clause is true''
(lines 35--36). \texttt{WidthAtMost k F} says every clause has at most
\(k\) literals as a set (line 38). Over a finite variable type,
\texttt{satisfyingAssignments F} filters the universe of all assignments
(lines 44--45), and \texttt{count F} is its cardinality (line 47): the
number of satisfying assignments among all \(2^n\) when
\(V = \mathtt{Fin}\ n\). The \texttt{Finset} choices change nothing in the
meaning: a repeated clause or literal collapses, the empty clause is
allowed and unsatisfiable, a tautological clause \(\{x, \neg x\}\) is
allowed and counts two toward the width. That is consistent with the
manuscript's ``Repeating a literal permits clauses with fewer than \(k\)
literals'' (p.~1, line~44).

\paragraph{The input encoding (lines 53--81).}
\texttt{encodeNat n} is \(n\) ones followed by a zero (line 55): unary.
\texttt{frame} doubles a bit string, prefixing each bit with a one, and
ends it with a zero (lines 57--58). \texttt{encodeList enc} prefixes each
element's frame with a one and ends the list with a zero (lines 60--62).
A literal becomes the number \(2\cdot\mathrm{var} + [\text{positive}]\),
at most \(2n - 1\) (lines 64--65); a clause is its sorted list of literal
numbers, a formula its sorted list of clauses (lines 67--71).
\texttt{encodeCNF F} is \texttt{encodeNat n} followed by the formula's
list of lists of unary numbers (lines 73--77), and
\texttt{formulaBitLength F} is the length of that bit string (line 79).
Two things follow. The number of variables \(n\) is part of the input,
so a variable that occurs in no clause is still counted. And the length
is polynomial: a clause of \(j \le k\) literals costs \(2j(4n+1)+4\) bits,
so
\[
  L \;=\; \mathtt{formulaBitLength}\ F \;\le\; n + 2 + m\,\bigl(2k(4n+1) + 4\bigr)
  \;=\; O(k\,n\,m)
\]
for \(m\) clauses. The solution itself proves
\((L+1)^d \le (8k+7)^d\,(n + |F| + 1)^{2d}\)
(\clonefile{SqrtLogSpace/Statement.lean}, lines 9--14). Emulating the
definitions (\repofile{reports/scratch/eth_audit/check_encoding.py},
random seed 1), one random formula with \(n = 125\), \(m = 500\),
\(k = 4\) encodes to \(L = 1{,}006{,}223\) bits and one with
\(n = 1{,}000\), \(m = 4{,}000\), \(k = 3\) to \(47.9\) million, both
under the bound; the length depends on the formula's literals, not on
\((n, m, k)\) alone. Nothing in the
encoding is exponential; the unary numbers are linear in \(n\).

\paragraph{The machine (lines 83--128).}
\texttt{PrimitiveMachine} (lines 92--99) is a finite object: a number
\texttt{tapeCount} of tapes and a number \texttt{stateCount} of states,
an input tape and a start state, a table \texttt{scan} saying which tape
each state reads, and a table \texttt{transition} from
\(\mathtt{Fin}\ \mathtt{stateCount} \times \mathtt{Option\ Bool}\)---a
domain of \(3 \cdot \mathtt{stateCount}\) points---into either a halting
answer (a \texttt{Bool}) or a next state with one action. The action type
is Mathlib's \texttt{Turing.TM0.Stmt (Option Bool)}: move the head left or
right, or write one of the three symbols \texttt{none}, \texttt{some
false}, \texttt{some true}, five values in all
(\path{Mathlib/Computability/TuringMachine/PostTuringMachine.lean}
lines 118--120 and \path{Tape.lean} lines 421--423 at Mathlib
\texttt{d13f23b7}, the clone's pinned revision). A configuration
\texttt{Cfg} (lines 103--105) is a state, halted or not, and one Mathlib
\texttt{Tape} per tape index (a head symbol and two blank-padded halves,
\path{Tape.lean} lines 408--414, blank \(=\) \texttt{none}).
\texttt{step} (lines 107--119) returns \texttt{none} in a halted state;
otherwise it reads the head of the one tape the state scans, looks the
pair up in the table, and either halts with the answer, or writes one
symbol or moves that one head by one cell and changes state. \texttt{init}
(lines 121--124) puts the input on the input tape with the head at its
first bit and every other tape blank (\texttt{Tape.mk}$_1$,
\path{Tape.lean} line 494). A Lean function on a finite domain is,
whatever term defines it, a finite table; so a step is one lookup in a
table of \(3 \cdot \mathtt{stateCount}\) rows, and no field of the
structure can carry unbounded computation.

\paragraph{Traces (lines 87--90).}
\texttt{Trace step t a c} is an inductive proposition with two
constructors: \texttt{refl}, a trace of length 0 from \(a\) to \(a\)
itself, and \texttt{cons}, which extends a trace from \(b\) by one
successful step \(a \mapsto b\). So it holds exactly when \(c\) is
reached from \(a\) by \(t\) successful applications of \texttt{step}, and
a trace of length 0 to a different configuration is impossible. Since
\texttt{step} is a function, the run is deterministic.

\paragraph{Space monitoring (lines 147--189).}
\texttt{LocatedCfg} (lines 149--153) pairs a configuration with, per tape,
the head's current position and the least and greatest positions seen so
far, as integers; \texttt{locatedInit} starts every head at 0 with the
input tape's range covering the input (lines 155--159);
\texttt{headDisplacement} is \(-1\), \(+1\) or \(0\) for a move left, a
move right or a write (lines 161--164); \texttt{monitor} (lines 166--182)
applies the same table lookup to update the position and the running
minimum and maximum of the scanned tape, and nothing else;
\texttt{locatedStep} is \texttt{step} on the configuration followed by
\texttt{monitor} (lines 184--186). \texttt{visitedSpace} (lines 188--189)
is the sum over tapes of \(\mathrm{upper} - \mathrm{lower} + 1\): the
number of cells whose head position was ever visited. \texttt{countWord a}
is the binary expansion of \(a\), least significant bit first, with
\texttt{[false]} for \(0\) (line 147), and \texttt{toNat} (lines 7--9)
decodes such a word.

\paragraph{The exponent and the bound (lines 130--141).}
\begin{quote}\small
\texttt{runtimeScale n = n * (Real.log (Real.log (n + 16)))\^{}2 /}\\
\texttt{\quad Real.sqrt (Real.log (n + 16))}
\end{quote}
(lines 132--134, transcribed; the casts to \(\mathbb{R}\) omitted), that is
\[
  \rhon \;=\; \frac{n\,\bigl(\ln\ln(n+16)\bigr)^2}{\sqrt{\ln(n+16)}},
\]
with \texttt{Real.log} the natural logarithm. The manuscript fixes the
same base: ``Throughout, log denotes the natural logarithm'' (p.~1, lines
48--49; again at p.~7, lines 326--327). The shift by 16 keeps
\(\ln\ln\) defined and positive at every \(n \ge 0\) (\(\ln 16 = 2.77\)),
so the theorem can be stated for all \(n\) with no proviso. Then
\begin{quote}
\texttt{runtimeBound C A d L n = C * (L + 1)\^{}d * 2\^{}(A * runtimeScale n)}
\end{quote}
(lines 136--139, transcribed), a real number, with the real-exponent
power.

\subsection{The theorem, with its quantifiers}
\label{app:tech-theorem}

This is the detail behind Section~\ref{sec:statement-theorem}. Lines
194--217 of the comparator file state, for every \(k\) with \(2 \le k\):
\[
\small
\begin{array}{l}
\exists\ \mathtt{machine},\ \mathtt{countTape},\ \mathtt{answerTape},\
  C, A, S \in \mathbb{R},\ d, e \in \mathbb{N},\\[2pt]
\quad 0 < C \ \wedge\ 0 < A \ \wedge\ 0 < S \ \wedge\ 0 < d \ \wedge\ 0 < e\ \wedge\\[2pt]
\quad \forall\, n \in \mathbb{N},\ \forall\, F : \mathtt{CNF}\ (\mathtt{Fin}\ n),\
  \mathtt{WidthAtMost}\ k\ F \;\rightarrow\\[2pt]
\qquad \exists\ \mathtt{time} \in \mathbb{N},\ \mathtt{final} : \mathtt{LocatedCfg}\ \mathtt{machine},\\[2pt]
\qquad\quad \mathtt{Trace}\ (\mathtt{locatedStep}\ \mathtt{machine})\ \mathtt{time}\
  (\mathtt{locatedInit}\ \mathtt{machine}\ (\mathtt{encodeCNF}\ F))\ \mathtt{final}\\[2pt]
\qquad\qquad\qquad\text{(lines 201--202)}\\[2pt]
\qquad\quad \wedge\ \text{the count tape holds \texttt{countWord (count F)}}
  \quad\text{(203--204)}\\[2pt]
\qquad\quad \wedge\ \text{the answer tape holds the bit \([0 < \mathtt{count}\ F]\)}
  \quad\text{(205--206)}\\[2pt]
\qquad\quad \wedge\ \text{the halting state carries that bit}
  \quad\text{(207)}\\[2pt]
\qquad\quad \wedge\ \mathtt{toNat}\ (\mathtt{countWord}\ (\mathtt{count}\ F)) = \mathtt{count}\ F
  \quad\text{(208)}\\[2pt]
\qquad\quad \wedge\ |\mathtt{countWord}\ (\mathtt{count}\ F)| \le n + 1
  \quad\text{(209)}\\[2pt]
\qquad\quad \wedge\ \mathtt{time} \le C\,(L+1)^d\,2^{A\,\rhon},\ L = \mathtt{formulaBitLength}\ F
  \quad\text{(210--211)}\\[2pt]
\qquad\quad \wedge\ \forall\ \text{prefix of the trace},\
  \mathtt{visitedSpace} \le S\,(L+1)^e
  \quad\text{(212--217).}
\end{array}
\]
The machine and the five constants are quantified after \(k\) and before
\(n\) and \(F\); inside the universal quantifier only \texttt{time} and
\texttt{final} are chosen. This is the standard \(O_k(1)\) reading: for
each \(k\) one fixed finite machine and one set of constants serve every
formula of every size. The table's size is \(O_k(1)\) in the same sense,
with no explicit function of \(k\). There is no lower bound on \(n\): the
manuscript's Theorem~1.1 and Proposition~6.4 are stated for \(n \ge 20\)
(p.~2, lines 56--57; p.~43, lines 2474--2475), and the solution's
auxiliary \texttt{exact\_count\_eth} keeps \(20 \le n\)
(\clonefile{SqrtLogSpace/Final.lean}, lines 104--109); the exported
theorem does not, and the \(+16\) in \(\rhon\) is what makes that
possible.

The proof's own constants, read for their values only: at the point where
the large-input run is assembled, the exponent constant is set to
\(A = 256 + A_t\) with \(A_t \ge 1\) taken from a leaf contract
(\clonefile{SqrtLogSpace/NativeLargeEntry.lean}, line 51;
\clonefile{SqrtLogSpace/NativeLeafContract.lean}, lines 792--793), and
\(1 \le C\), \(1 \le A\) are in the large-input run's own statement
(\clonefile{SqrtLogSpace/NativeLargeEntry.lean}, line 12). The witness
machine is compiled from a program over 80 one-bit registers and 1,220
tapes, with \(\mathtt{countTape} = 256\) and \(\mathtt{answerTape} = 1216\)
(\clonefile{SqrtLogSpace/Final.lean}: the program's type at line 11, the
machine built from it at lines 98--101); its state count is
a finite type's cardinality that nothing computes. The solution's
internal \(\rho\) (\clonefile{SqrtLogSpace/Asymptotics.lean}, lines
10--11) is in base 2 and is related to \texttt{runtimeScale} by
\(\mathtt{runtimeScale}\ n \le 4\,\rho_2(n)\) for \(n \ge 20\)
(lines 190--191); the base changes the exponent by a constant factor
that \(A\) absorbs.

The clause-count bound used in Sections~\ref{sec:statement} and
\ref{sec:consequences}: a \texttt{Finset} of clauses of width at most
\(k\) over \(n\) variables has at most \(\sum_{j \le k}\binom{2n}{j} =
O(n^k)\) members (\(2{,}604{,}376\) at \(n = 125\), \(k = 3\)), so
\(m = O(n^k)\), \(L = O(k\,n\,m) = O(n^{k+1})\), and the whole bound is
\(n^{O_k(1)}\,2^{A\rhon}\), which is \(2^{o(n)}\) in \(n\) alone because
\(\rhon/n \to 0\) \leanref{rho-little-o}.

\subsection{The spine of the proof, with every line}
\label{app:tech-spine}

This is the detail behind Section~\ref{sec:proof-spine}. Unqualified
paths are relative to \path{lean/OAI/Algorithms/SubexponentialSAT/} in
the clone.

\paragraph{The machine.} The theorem's machine is
\texttt{GuardedCompiler.primitive program} with the count tape
\texttt{256} and the answer tape \texttt{1216} (both \texttt{: Fin 1220};
\clonefile{SqrtLogSpace/Final.lean}, line~100), where \texttt{program}
is \texttt{Native.countingProgram cutoff k} \texttt{(NativeLeafProgram.command k)}
(line~40), a \texttt{BitProgram.Program 80 1220 labels} (line~11): a
first-order program over 80 one-bit registers, 1,220 Boolean stacks and
a finite label set, with twelve instruction kinds
(\clonefile{BitProgram.lean}, lines 10--26). \texttt{primitive}
(\clonefile{SqrtLogSpace/GuardedCompiler.lean}, line~131) is a
\texttt{noncomputable abbrev}: it takes the transition table
\texttt{transition p} (lines 30--47, one case per instruction) and the
stack selector \texttt{selected p} (lines 14--17), feeds them to the
stack compiler's controller, and renames the finite state type to
\texttt{Fin stateCount} through \texttt{Fintype.equivFin} (line~132).
The machine is thus a concrete table defined by pattern matching on the
program, composed with a classical renaming of states; the theorem is
existential, so the noncomputability costs it nothing.

\paragraph{The time bound, and the constants.} For \texttt{cutoff < n},
\path{NativeLargeEntry.run} (\clonefile{SqrtLogSpace/NativeLargeEntry.lean},
lines 11--33) supplies a trace of length at most
\texttt{runtimeBound C A d L n} (line~27) with $1 \le C \wedge 1 \le A$
(line~12). Its exponent constant is \texttt{256 + At} (line~51), where
\texttt{At} comes from the leaf program's uniform contract
(\path{NativeLeafContract.uniform_contract}, line~39; the contract
gives $1 \le \mathtt{At}$, \clonefile{SqrtLogSpace/NativeLeafContract.lean},
line~793), and its cutoff is \texttt{max 20 (max Nleaf Ns)} (line~50),
with \texttt{Ns} obtained non-constructively from an eventual inequality
(\path{Filter.eventually_atTop}, line~48). The $2^{A\cdot\mathtt{runtimeScale}\,n}$
factor is assembled by \path{runtime_of_tree_and_leaf} (line~46)
from the DFS visit bound of \path{Native.outer_from_setup}
(\clonefile{SqrtLogSpace/NativeOuter.lean}, lines 53--56: visits at most
a square polynomial in $n$ times $2^{n/\mathtt{sparsifySchedule}\,n}$)
and the leaf time. For $n \le \mathtt{cutoff}$,
\path{Native.entry_small_primitive}
(\clonefile{SqrtLogSpace/NativeEntry.lean}, lines 81--100) gives
$\mathtt{time} \le 28700\,(L+1)^2$ (line~92) through the finite scanner
of \clonefile{SqrtLogSpace/FiniteFallback.lean}. The two branches are
folded by \texttt{runtime\_with\_fallback} and
\texttt{space\_with\_fallback}
(\clonefile{SqrtLogSpace/NativeOuterBudgets.lean}, lines 297 and 317)
into one triple, and the witnesses the final theorem exhibits are
\texttt{C + 28700, A + 1, S + 5022741, d + 2, e + 2}
(\clonefile{SqrtLogSpace/Final.lean}, line~41). So the proof as written
has an exponent constant of at least $258$ in front of $\rhon$, a
polynomial degree \texttt{d} that depends on $k$ (line~51: $\mathtt{db} + (4(k-2)+2) + \mathtt{df} + 2$) and a leading constant of at least
$28{,}701$; none of these is in the statement, which asserts only that
they are positive, and none of \texttt{At}, \texttt{db}, \texttt{df} is
computed anywhere.

\paragraph{Where a hidden cost model would live, and why it cannot.}
The compile lemmas conclude in
the comparator's own vocabulary: \path{GuardedCompiler.primitive_execution}
(\clonefile{SqrtLogSpace/GuardedCompiler.lean}, lines 160--162) turns a
source run of \texttt{n} steps into \texttt{RunsIn (primitive p).step
(7 * n)}, with the overhead \texttt{7} proved as a theorem
(\texttt{compiler\_overhead}, lines 152--153), and
\path{CompileSpace.compile_word_answer_space}
(\clonefile{SqrtLogSpace/CompileSpace.lean}, lines 125--143) produces a
\texttt{Computation.Trace (locatedStep (primitive p)) time \ldots} with
$\mathtt{time} \le 7\cdot\mathtt{sourceSteps}$ and
$\mathtt{visitedSpace}\ \mathtt{intermediate} \le s\,(\mathtt{bound} + \mathtt{input.length} + 15)$ on
every prefix (lines 134, 135 and 143); \path{primitive_located_step}
(\clonefile{SqrtLogSpace/GuardedCompiler.lean}, line~600) is the per-step
form. \texttt{Trace}, \texttt{locatedStep} and \texttt{visitedSpace} are
the challenge's definitions (its lines 87--90 and 149--189), and
\texttt{time} in the final theorem is literally the length of a
\texttt{Computation.Trace} of \texttt{PrimitiveMachine.step} (its
line~201). No cost function of the compiler's appears in the statement;
a wrong cost model would make these lemmas unprovable, not the theorem
false.

\paragraph{Correctness of the count.} \path{Native.outer_from_setup}
(\clonefile{SqrtLogSpace/NativeOuter.lean}, lines 50--51) yields
\texttt{final.stacks ABI.countStack = countWord (SAT.count F)} and
\texttt{final.stacks ABI.answerStack = [decide (0 < SAT.count F)]} from
the canonical DFS execution, and
\path{Native.counting_primitive_from_entry} (line~73) transports
both to the primitive machine's tapes; \texttt{countWord\_toNat} and
\texttt{countWord\_count\_length} supply the decoding and length clauses
of the statement.

\paragraph{The statement match, counted.} The solution file does not
import the challenge file: \clonefile{SqrtLogSpace/Final.lean} begins
(lines 1--3) with \texttt{NativeLargeEntry}, \texttt{SqrtLogSpace.Statement}
and \texttt{SqrtLog.StatementCanonical}. The library re-declares every
definition the statement uses under the same fully qualified names, in
seven files: the SAT vocabulary in \clonefile{SAT.lean}, the bit
arithmetic in \clonefile{BinaryArithmetic.lean}, the nine encoding
functions in \clonefile{Encoding.lean}, the machine in
\clonefile{Computation.lean}, \texttt{runtimeScale} and
\texttt{runtimeBound} in \clonefile{SqrtLog/StatementScale.lean}, the
space instrumentation in \clonefile{SqrtLogSpace/VisitedTape.lean} and
\texttt{countWord} in \clonefile{SqrtLogSpace/EncodingBounds.lean}
(line~17, the challenge's line~147). The definition diff compares 35
declarations: 33 are identical character for character after whitespace
normalisation, and the remaining two, \texttt{runtimeScale} and
\texttt{runtimeBound}, differ only in that the challenge writes
\texttt{noncomputable def} where the library opens a
\texttt{noncomputable section} (\clonefile{SqrtLog/StatementScale.lean},
line~15). The theorem statement itself, from \texttt{theorem
exact\_count\_and\_sat} to \texttt{:= by}, is identical in the two files,
1,353 characters raw and 1,124 after normalisation, at line~194 of the
challenge and line~74 of \clonefile{SqrtLogSpace/Final.lean}; the
challenge's proof is the single word \texttt{sorry} at its line~218.

\subsection{The keyword sweep, item by item}
\label{app:tech-keywords}

This is the detail behind Section~\ref{sec:proof-keywords}. The only
\texttt{set\_option}s in the closure are two
\texttt{set\_option quotPrecheck false}
(\clonefile{SqrtLog/LeafProgramPreparedRecoveryQuery.lean}, line~15, and
\clonefile{SqrtLog/LeafProgramPreparedRecoveryWitness.lean}, line~14),
which switch off a pre-check on local notation and have no bearing on
soundness; there is no \texttt{maxHeartbeats}, \texttt{maxRecDepth} or
\texttt{autoImplicit} option anywhere in the closure, and the lakefile
sets \texttt{autoImplicit} to \texttt{false} for the whole package
(\clonefile{lean/lakefile.lean}, line~9). Nine declarations are
\texttt{opaque}, and every one carries a body: three recovery functions
(\texttt{queryRecovery} at line~59 of
\clonefile{SqrtLog/LeafProgramPreparedRecoveryQuery.lean}, with its
\texttt{:= by} at line~86; \texttt{rawRecovery} and
\texttt{fixedRecovery} at lines 178 and 198 of
\clonefile{SqrtLog/LeafProgramPreparedRecoveryWitness.lean}) and six
\texttt{private opaque} constants of a subtype
\texttt{\{ports // ports = \ldots\}}, each with an explicit value
(line~22 of \clonefile{SqrtLog/LeafProgramQueryCleanup.lean} and five
others). An \texttt{opaque} with a body is a definition whose unfolding
is hidden from elaboration; the kernel still checks the body against the
type. The 1,048 \texttt{noncomputable} declarations, the eight
\texttt{open Classical}, and the four explicit uses each of
\texttt{Classical.choice} and \texttt{propext} are the classical
reasoning the comparator permits. The 148 \texttt{instance}s and 40
\texttt{local instance}s are \texttt{Fintype}, \texttt{Decidable} and
algebraic instances; an instance cannot be ``inconsistent'' without an
axiom, and there is none. The 14,034 lines with \texttt{decide} (as \repofile{closure.py} counts them, comments stripped) are
kernel-evaluated decision procedures, not the native evaluator. Nine
\texttt{termination\_by} and five \texttt{decreasing\_by} clauses are the
recursions the kernel checks.

The lakefile's two hooks: a \texttt{run\_cmd} block (lines 290--304) that
evaluates dependency declarations with \texttt{unsafe env.evalConstCheck}
and calls \path{preparePatchedDependencies}, which clones and patches
eleven dependencies, and a \texttt{post\_update} hook (line~306 onward)
that verifies or applies the patches for the rest. They touch only
\texttt{.lake/packages}, and only non-Mathlib packages. The 23
compatibility patches under \clonefile{lean/patches/} are for third-party
packages (\texttt{AINTLIB}, \texttt{PrimeNumberTheoremAnd},
\texttt{carleson}, \texttt{iut}, \ldots) that this closure does not
import. The registry \clonefile{lean/formalization.yaml} (lines
1683--1685) points at the same declaration and file as the comparator's
JSON.

\subsection{The hypotheses in quantifier form}
\label{app:tech-hypotheses}

This is the detail behind Section~\ref{sec:consequences-hypotheses}.
The property used throughout: for each fixed $k \ge 2$ there is one
deterministic machine which, on the canonical encoding of any
\texttt{WidthAtMost}~$k$ formula $F$ on $n$ variables (comparator line
38), halts with \texttt{countWord (count F)} on its count tape (lines
203--204, 147), the bit $[0 < \mathtt{count}\ F]$ on its answer tape and
in its halting state (lines 205--207), within $C\,(L+1)^d\,2^{A\rhon}$
steps (lines 210--211). Since $\rhon/n \to 0$ \leanref{rho-little-o},
for every $\delta > 0$ there is an $n_0(k, \delta)$ beyond which the
bound is below $2^{\delta n}$; below $n_0$ the bound is a constant
depending on $k$ and $\delta$. So the running time is
$O_{k,\delta}(2^{\delta n})$ for every $\delta > 0$, and
$2^{\delta n}(n+m)^{O_k(1)}$ with the polynomial written out.

\paragraph{ETH.} In quantifiers: there is a $\delta > 0$ such that no
deterministic algorithm decides 3-SAT within $2^{\delta n}\,\mathrm{poly}(n+m)$
steps on all inputs. Its negation: for every $\delta > 0$ there is one.
The verified statement at $k = 3$ supplies, for every $\delta > 0$, the
same machine, whose answer bit is the decision, within
$O_{\delta}(2^{\delta n})\,\mathrm{poly}(n+m)$ steps. This holds on every
input of every size, so it refutes the hypothesis under either reading
of ``$2^{\Omega(n)}$'' (a lower bound at infinitely many $n$, or at all
large $n$). In the vocabulary of \citet{ImpagliazzoPaturi2001}, where
$s_k$ is the infimum of the $\delta$ for which $k$-SAT has an
$O(2^{\delta n})$ algorithm and ETH is $s_3 > 0$: the verified statement
gives $s_k = 0$ for every $k \ge 2$, by one algorithm per $k$.

\paragraph{SETH.} The manuscript's form: for every
$\varepsilon \in (0,1)$, there is a fixed $k$ for which no deterministic
algorithm solves $k$-SAT in $2^{(1-\varepsilon)n}(n+m)^{O(1)}$ time. The
negation: there is an $\varepsilon \in (0,1)$ such that for every fixed
$k$ there is a deterministic $2^{(1-\varepsilon)n}(n+m)^{O(1)}$
algorithm. The verified statement gives more: for every
$\varepsilon \in (0,1)$ and every $k \ge 2$, the machine for $k$ runs
within $C_k(L+1)^{d_k}\,2^{A_k\rhon}$ steps, and
$A_k\rhon \le (1-\varepsilon)n$ for all $n \ge n_0(k, \varepsilon)$
because $\rhon = o(n)$; below $n_0$ the exponential factor is a
constant. So each $k$ has a $2^{(1-\varepsilon)n}(n+m)^{O_k(1)}$
algorithm, and the polynomial's degree depending on $k$ is exactly what
the hypothesis's ``for a fixed $k$'' permits. For the record, SETH
implies ETH: if 3-SAT had a $2^{o(n)}$ algorithm then, by the
Sparsification Lemma of \citet{ImpagliazzoPaturiZane2001}, which reduces
$k$-SAT on $n$ variables to $2^{\varepsilon n}$ instances with
$O_{k,\varepsilon}(n)$ clauses and so, after the clause-splitting
reduction, to 3-SAT on $O_{k,\varepsilon}(n)$ variables, so would
$k$-SAT for every fixed $k$.

\paragraph{Gap-ETH.} The negation of the manuscript's form: for every
$\delta, \gamma > 0$ there is an $O(2^{\gamma n})$ algorithm that
separates a satisfiable 3-CNF from one in which every assignment
violates at least a $\delta$ fraction of the clauses, already at
$m = O(n)$. Exact decision separates them, since a formula in which
every assignment violates a positive fraction of the clauses is
unsatisfiable; and the verified machine at $k = 3$ decides within
$n^{O(1)}\,2^{A\rhon} = O_{\gamma}(2^{\gamma n})$ steps for every
$\gamma > 0$, at every $m$ (at $m = O(n)$ the polynomial is
$O(n^{2d})$). The promise is never used: the machine answers correctly
on every input, promised or not, with probability $1$, which refutes
the randomised form as well.

\paragraph{NSETH.} The negation of the manuscript's form: there is an
$\varepsilon > 0$ such that every fixed $k$ has a nondeterministic
$2^{(1-\varepsilon)n}(n+m)^{O(1)}$ algorithm for unsatisfiability. The
verified machine's answer bit is \texttt{false} exactly when
$\mathtt{count}\ F = 0$ (lines 205--207), so it decides unsatisfiability
deterministically within the SETH bound above for every $\varepsilon$
and every $k$; and a deterministic machine is a nondeterministic one
whose only computation path is the deterministic run, so every
acceptance convention for nondeterministic time is met.

\paragraph{\#ETH.} The negation: for every $\delta > 0$ there is a
deterministic $O(2^{\delta n})$ algorithm that outputs the exact count.
The verified machine at $k = 3$ halts with \texttt{countWord (count F)}
on its count tape, a binary word of length at most $n + 1$ that decodes
to $\mathtt{count}\ F$ (lines 203--204, 208--209), within the same bound.

\paragraph{Shape (S) in detail.} For a theorem ``SETH $\Rightarrow$ no
$(c-\varepsilon)^k$ algorithm for $Y$'', a $(c-\varepsilon)^k$ algorithm
for $Y$ obtained from the verified statement would need the exponent of
the SAT instance to be below $k \log_2 c$ by a constant factor, and the
verified statement's exponent is $A\rho(n')$ in the variable count $n'$
of the SAT instance with $A$ unknown; no value of $A$ is on record below
$257$, and no base $c$ is derived from the verified statement for any
$Y$.

\subsection{How the two guarantee curves are scaled}
\label{app:tech-scaling}

This is the detail behind Section~\ref{sec:practitioner-guarantees}.
Both guarantees are taken at constant 1 (the polynomial factors, $C$,
$A$ and $d$ set aside) and both are multiplied by the same factor
$t_{20}/2^{20}$, where $t_{20} = 1.28 \times 10^{-4}$~s is CaDiCaL's
median at $n = 20$. The factor is one number for both curves, so their
ratio is the ratio of the guarantees themselves, $2^{\rhon - n}$: below
1 for $n < 50$, just above 1 at $n = 50$ ($\rho(50) = 50.14$, ratio
$1.10$), and above 1 from there ($2.8 \times 10^{3}$ at
$n = 100$ and $4.0 \times 10^{51}$ at $n = 500$). A reader who prefers
each curve scaled to the $n = 20$ point on its own would multiply
$2^{\rhon}$ by a further $2^{20 - \rho(20)} = 2^{2.79}$ and see the two
curves touch at $n = 20$ and the verified one above the trivial one
from there on; the common factor keeps the crossing where the statement
puts it (\texttt{ETH\_AMBIGUOUS.MD}, 06a). In that scale the
brute-force guarantee reaches one second near $n = 33$ and the verified
guarantee near $n = 35$; at $n = 60$ the measured CaDiCaL median is
$3 \times 10^{-4}$~s against $2^n$ at $1.4 \times 10^{8}$~s and
$2^{\rhon}$ at $5.4 \times 10^{8}$~s, 11.6 and 12.2 decades
respectively; at $n = 500$ the decision median is above the budget
(8 of 20 finished within 60~s), against $10^{140.6}$~s and
$10^{192.2}$~s. Regenerate: \texttt{python3 -I paper3/scripts/practitioner.py}.

\par\endgroup
